The numerical data, source code, case configurations, and supporting documentation for this study are publicly available in the \texttt{NewtonPicard} repository on GitHub (\url{https://github.com/voronoimeshmaker/NewtonPicard}). The versioned archive is available on Zenodo as version~1.0.0 (DOI: \href{https://doi.org/10.5281/zenodo.23306169}{10.5281/zenodo.23306169}), corresponding to Git commit \href{https://github.com/voronoimeshmaker/NewtonPicard/commit/cc2ae381c79f10cf8f131d4b4906a6abc202881b}{\texttt{cc2ae381c79f}}. The archived results include the numerical campaigns performed using ADOL-C for source-term differentiation. The computational framework uses \texttt{VoronoiMeshMaker} for Vo\-ro\-noi mesh generation with CGAL~\cite{TheCGALProject2024CGALManual}, \texttt{VoroGradToolKit} for gradient reconstruction, and \texttt{GSMakerLib} for assembling and solving the discrete Grad--Shafranov systems. PETSc, SLEPc, and ADOL-C provide the numerical linear algebra, spectral analysis, and automatic differentiation facilities, respectively.
